% piano di progetto
% da compilare con il comando pdflatex Piano_di_progetto_v_x.x.x.tex

% Dichiarazioni di ambiente e inclusione di pacchetti
% da usare tramite il comando % Dichiarazioni di ambiente e inclusione di pacchetti
% da usare tramite il comando % Dichiarazioni di ambiente e inclusione di pacchetti
% da usare tramite il comando \input{../../util/hx-ambiente}

\documentclass[a4paper,titlepage]{article}
\usepackage[T1]{fontenc}
\usepackage[utf8]{inputenc}
\usepackage[english,italian]{babel}
\usepackage{microtype}
\usepackage{lmodern}
\usepackage{underscore}
\usepackage{graphicx}
\usepackage{eurosym}
\usepackage{float}
\usepackage{fancyhdr}
\usepackage[table]{xcolor}
\usepackage{longtable}
\usepackage{chngpage}
\usepackage{grffile}
\usepackage[titles]{tocloft}
\usepackage{hyperref}
\hypersetup{hidelinks}

\usepackage{../../util/hx-vers}
\usepackage{../../util/hx-macro}
\usepackage{../../util/hx-front}

% solo se si vuole una nuova pagina ad ogni \section:
\usepackage{titlesec}
\newcommand{\sectionbreak}{\clearpage}

% stile di pagina:
\pagestyle{fancy}

% solo se si vuole eliminare l'indentazione ad ogni paragrafo:
\setlength{\parindent}{0pt}

% intestazione:
\lhead{\Large{\proj}}
\rhead{\includegraphics[keepaspectratio=true,width=50px]{../../util/hivex_logo2.png}}
\renewcommand{\headrulewidth}{0.4pt}

% pie' di pagina:
\lfoot{\email}
\rfoot{\thepage}
\cfoot{}
\renewcommand{\footrulewidth}{0.4pt}

% spazio verticale tra le celle di una tabella:
\renewcommand{\arraystretch}{1.5}

% profondità di indicizzazione:
\setcounter{tocdepth}{4}
\setcounter{secnumdepth}{4}


\documentclass[a4paper,titlepage]{article}
\usepackage[T1]{fontenc}
\usepackage[utf8]{inputenc}
\usepackage[english,italian]{babel}
\usepackage{microtype}
\usepackage{lmodern}
\usepackage{underscore}
\usepackage{graphicx}
\usepackage{eurosym}
\usepackage{float}
\usepackage{fancyhdr}
\usepackage[table]{xcolor}
\usepackage{longtable}
\usepackage{chngpage}
\usepackage{grffile}
\usepackage[titles]{tocloft}
\usepackage{hyperref}
\hypersetup{hidelinks}

\usepackage{../../util/hx-vers}
\usepackage{../../util/hx-macro}
\usepackage{../../util/hx-front}

% solo se si vuole una nuova pagina ad ogni \section:
\usepackage{titlesec}
\newcommand{\sectionbreak}{\clearpage}

% stile di pagina:
\pagestyle{fancy}

% solo se si vuole eliminare l'indentazione ad ogni paragrafo:
\setlength{\parindent}{0pt}

% intestazione:
\lhead{\Large{\proj}}
\rhead{\includegraphics[keepaspectratio=true,width=50px]{../../util/hivex_logo2.png}}
\renewcommand{\headrulewidth}{0.4pt}

% pie' di pagina:
\lfoot{\email}
\rfoot{\thepage}
\cfoot{}
\renewcommand{\footrulewidth}{0.4pt}

% spazio verticale tra le celle di una tabella:
\renewcommand{\arraystretch}{1.5}

% profondità di indicizzazione:
\setcounter{tocdepth}{4}
\setcounter{secnumdepth}{4}


\documentclass[a4paper,titlepage]{article}
\usepackage[T1]{fontenc}
\usepackage[utf8]{inputenc}
\usepackage[english,italian]{babel}
\usepackage{microtype}
\usepackage{lmodern}
\usepackage{underscore}
\usepackage{graphicx}
\usepackage{eurosym}
\usepackage{float}
\usepackage{fancyhdr}
\usepackage[table]{xcolor}
\usepackage{longtable}
\usepackage{chngpage}
\usepackage{grffile}
\usepackage[titles]{tocloft}
\usepackage{hyperref}
\hypersetup{hidelinks}

\usepackage{../../util/hx-vers}
\usepackage{../../util/hx-macro}
\usepackage{../../util/hx-front}

% solo se si vuole una nuova pagina ad ogni \section:
\usepackage{titlesec}
\newcommand{\sectionbreak}{\clearpage}

% stile di pagina:
\pagestyle{fancy}

% solo se si vuole eliminare l'indentazione ad ogni paragrafo:
\setlength{\parindent}{0pt}

% intestazione:
\lhead{\Large{\proj}}
\rhead{\includegraphics[keepaspectratio=true,width=50px]{../../util/hivex_logo2.png}}
\renewcommand{\headrulewidth}{0.4pt}

% pie' di pagina:
\lfoot{\email}
\rfoot{\thepage}
\cfoot{}
\renewcommand{\footrulewidth}{0.4pt}

% spazio verticale tra le celle di una tabella:
\renewcommand{\arraystretch}{1.5}

% profondità di indicizzazione:
\setcounter{tocdepth}{4}
\setcounter{secnumdepth}{4}

\version{0.0.1}
\author{\LB, \PB}
\supervisor{\GG, \MM}
\dest{Uso esterno}
\title{Piano di progetto}

\renewcommand{\arraystretch}{1.5}
\setcounter{tocdepth}{4}
\setcounter{secnumdepth}{4}

\begin{document}
\maketitle
% diario delle modifiche per il piano di progetto
% da includere con % diario delle modifiche per il piano di progetto
% da includere con % diario delle modifiche per il piano di progetto
% da includere con \include{diario}

\begin{diario}
	1.0.1 & {\GG} (Responsabile) & 03/02/2017 & Correzione degli errori individuati nelle prime tre sezioni del documento \\ \hline
	1.0.0 & {\PB} (Responsabile) & 09/01/2017 & Approvazione documento \\ \hline
	0.1.0 & {\MM} (Verificatore) & 07/01/2017 & Verifica documento \\ \hline
	0.0.9 & {\PB} (Responsabile) & 05/01/2017 & Stesura Organigramma \\ \hline
	0.0.8 & {\LB} (Responsabile) & 05/01/2017 & Stesura Consuntivo di Periodo \\ \hline
	0.0.7 & {\LB} (Responsabile) & 04/01/2017 & Stesura Preventivo di Periodo \\ \hline
	0.0.6 & {\LB} (Responsabile) & 03/01/2017 & Inserimento Gantt Diagrammi delle Attività \\ \hline
	0.0.5 & {\PB} (Responsabile) & 29/12/2016 & Stesura Pianificazione delle Attività \\ \hline
	0.0.4 & {\PB} (Responsabile) & 28/12/2016 & Stesura Introduzione \\ \hline
	0.0.3 & {\LB} (Responsabile) & 27/12/2016 & Stesura Modello di sviluppo \\ \hline
	0.0.2 & {\PB} (Responsabile) & 27/12/2016 & Stesura Analisi dei Rischi \\ \hline
	0.0.1 & {\LB} (Responsabile) & 26/12/2016 & Stesura scheletro \\ \hline
\end{diario}

\begin{diario}
	1.0.1 & {\GG} (Responsabile) & 03/02/2017 & Correzione degli errori individuati nelle prime tre sezioni del documento \\ \hline
	1.0.0 & {\PB} (Responsabile) & 09/01/2017 & Approvazione documento \\ \hline
	0.1.0 & {\MM} (Verificatore) & 07/01/2017 & Verifica documento \\ \hline
	0.0.9 & {\PB} (Responsabile) & 05/01/2017 & Stesura Organigramma \\ \hline
	0.0.8 & {\LB} (Responsabile) & 05/01/2017 & Stesura Consuntivo di Periodo \\ \hline
	0.0.7 & {\LB} (Responsabile) & 04/01/2017 & Stesura Preventivo di Periodo \\ \hline
	0.0.6 & {\LB} (Responsabile) & 03/01/2017 & Inserimento Gantt Diagrammi delle Attività \\ \hline
	0.0.5 & {\PB} (Responsabile) & 29/12/2016 & Stesura Pianificazione delle Attività \\ \hline
	0.0.4 & {\PB} (Responsabile) & 28/12/2016 & Stesura Introduzione \\ \hline
	0.0.3 & {\LB} (Responsabile) & 27/12/2016 & Stesura Modello di sviluppo \\ \hline
	0.0.2 & {\PB} (Responsabile) & 27/12/2016 & Stesura Analisi dei Rischi \\ \hline
	0.0.1 & {\LB} (Responsabile) & 26/12/2016 & Stesura scheletro \\ \hline
\end{diario}

\begin{diario}
	1.0.1 & {\GG} (Responsabile) & 03/02/2017 & Correzione degli errori individuati nelle prime tre sezioni del documento \\ \hline
	1.0.0 & {\PB} (Responsabile) & 09/01/2017 & Approvazione documento \\ \hline
	0.1.0 & {\MM} (Verificatore) & 07/01/2017 & Verifica documento \\ \hline
	0.0.9 & {\PB} (Responsabile) & 05/01/2017 & Stesura Organigramma \\ \hline
	0.0.8 & {\LB} (Responsabile) & 05/01/2017 & Stesura Consuntivo di Periodo \\ \hline
	0.0.7 & {\LB} (Responsabile) & 04/01/2017 & Stesura Preventivo di Periodo \\ \hline
	0.0.6 & {\LB} (Responsabile) & 03/01/2017 & Inserimento Gantt Diagrammi delle Attività \\ \hline
	0.0.5 & {\PB} (Responsabile) & 29/12/2016 & Stesura Pianificazione delle Attività \\ \hline
	0.0.4 & {\PB} (Responsabile) & 28/12/2016 & Stesura Introduzione \\ \hline
	0.0.3 & {\LB} (Responsabile) & 27/12/2016 & Stesura Modello di sviluppo \\ \hline
	0.0.2 & {\PB} (Responsabile) & 27/12/2016 & Stesura Analisi dei Rischi \\ \hline
	0.0.1 & {\LB} (Responsabile) & 26/12/2016 & Stesura scheletro \\ \hline
\end{diario}
\tableofcontents

\section{Introduzione}
	\subsection{Scopo del documento}
	Il presente documento illustra la pianificazione adottata dal gruppo {\hx} per la produzione del progetto {\proj}. [more?] In esso vi sono contenuti:
\begin{itemize}
	\item Scelta del modello di sviluppo adottato;
	\item Analisi dei rischi;
	\item Pianificazione delle attività;
	\item Preventivo [blablabla];
    \item Consuntivo di periodo.
\end{itemize}

	\subsection{Scopo del prodotto}
	\scopo
	\subsection{Glossario}
	[idem]
	\subsection{Riferimenti}
		\subsubsection{Normativi}
		[idem]
		\subsubsection{Informativi}


	\subsection{Scadenze}
	[idem]

\section{Modello di sviluppo}
[si sceglie incrementale]


\section{Analisi dei rischi}
	\subsection{A}
	\subsection{B}
		\subsubsection{B1}
		\subsubsection{B2}
	\subsection{C}

\section{Pianificazione}
[Si rimanda a Norme di progetto v_1_0_0 per il dettaglio del funzionamento di Asana]
\section{Preventivo}
\section{Consuntivo di periodo}

\end{document}
